\section{Calculus on \(\RR^n\)}

Consider \(\functiondomain{f}{\RR^n}{\RR}\), and we write
\[
    \vb{x} = \pqty{x_1, x_2, \ldots, x_n} \in \RR^n.
\]

Assume \(f\) is sufficiently differentiable (and in particular, at least \(\contClass{2}\), meaning \(f\) is twice continuously differentiable).

\begin{definition}[Stationary Point]
    \label{def:stationary-point}%
    We say a point \(\vb{a} \in \RR^n\) is a \emph{stationary point} of \(f\), if
    \[
        \pqty{\grad f} \pqty{\vb{a}} = \vb{0},
    \]
    where the \emph{gradient} is given by
    \[
        \grad f = \pqty{\pdv{f}{x_1}, \pdv{f}{x_2}, \ldots, \pdv{f}{x_n}}.
    \]
\end{definition}

We may also Taylor expand a function about a point, having
\[
    f\pqty{\vb{x}} = f\pqty{\vb{a}} + \sum_{i} \pqty{x_i - a_i} \pqty{\nabla_i f} \pqty{\vb{a}} + \frac{1}{2} \sum_{i, j} \pqty{x_i - a_i} \pqty{x_j - a_j} H_{ij} \pqty{\vb{a}} + \bigO \pqty{\norm{\vb{x} - \vb{a}}^3},
\]
where the \emph{Hessian matrix} is given by
\[
    H_{ij} = \pdv{f}{x_i}{x_j}.
\]

Since \(f\) is \(\contClass{2}\), \(H_{ij} = H_{ji}\). From now on, we simplify the notation with \(\vb{a} = \vb{0}\) by a translation to the origin. Under the summation convention, at a stationary point, we have
\[
    f\pqty{\vb{x}} - f\pqty{\vb{0}} = \frac{1}{2} x_i H_{ij} x_j + \bigO \pqty{\abs{\vb{x}}^3}.
\]

Recall that we may diagonalise \(H_{ij}\), since it is symmetric, and we may rotate the coordinates under the transformation
\[
    x_i = R_{ij} x'_{j}.
\]

Then the new Hessian matrix will be transformed as
\[
    \vb{H}' = \vb{R}\transpose \vb{H} \vb{R} = \pmqty{\dmat{\lambda_1, \lambda_2, \ddots, \lambda_n}}
\]
and the eigenvalues \(\lambda_1, \ldots, \lambda_n\) are the eigenvalues of \(\vb{H}\). Note that since \(\vb{H}\) is real, symmetric, all the eigenvalues are therefore real.

After diagonalisation, the quadratic form will be diagonalised as
\[
    f\pqty{\vb{x}} - f\pqty{\vb{0}} = \frac{1}{2} \sum_{i = 1}^{n} \lambda_i x_i'^2 + \bigO \pqty{\abs{\vb{x}}^3}.
\]

Therefore, we have the following result.

\begin{proposition}[Classification of Stationary Point using Hessian Matrix]
    \label{prop:classification-stationary-point}%
    Using the Hessian matrix, at a stationary point,
    \begin{itemize}
        \item if \(\lambda_i > 0\) for all \(i\), then \(f\pqty{\vb{x}} > f\pqty{\vb{0}}\) for sufficiently small \(\vb{x}\), and hence \(\vb{x} = \vb{0}\) is a \emph{local minimum};
        \item if \(\lambda_i < 0\) for all \(i\), then \(\vb{x} = \vb{0}\) is a \emph{local maximum};
        \item if there is some \(\lambda_i > 0\) and some other \(\lambda_j < 0\), then \(\vb{x} = \vb{0}\) is a \emph{saddle point}.
    \end{itemize}
\end{proposition}

\begin{remark}
    The remaining cases are \emph{degenerate} stationary points -- behaviour depends on \(\bigO\pqty{x^3}\) terms.
\end{remark}

\begin{corollary}[Hessian Matrix in Two-Variable Case]
    \label{cor:hessian-2-d}%
    In the special case where \(n = 2\), we have
    \begin{itemize}
        \item if \(\det \vb{H} > 0\) and \(\tr \vb{H} > 0\), then it is a local minimum;
        \item if \(\det \vb{H} > 0\) and \(\tr \vb{H} < 0\), then it is a local maximum;
        \item if \(\det \vb{H} < 0\), then it is a saddle point;
        \item if \(\det \vb{H} = 0\), then it is degenerate.
    \end{itemize}
\end{corollary}

\begin{proof}
    We have \(\det \vb{H} = \lambda_1 \lambda_2\), and \(\tr \vb{H} = \lambda_1 + \lambda_2\). Applying \autoref{prop:classification-stationary-point} gives the result.
\end{proof}

\begin{remark}
    If a local minimum is not a global minimum, then either:
    \begin{itemize}
        \item there is another stationary point that is the global minimum; or
        \item the global minimum occurs on the boundary of the domain of \(f\); or
        \item \(f\) does not have a global minimum.
    \end{itemize}
\end{remark}

\begin{example}[Finding and Classifying Stationary Points]
    \label{ex:finding-classifying-stationary-points}%
    Consider the function \(\functiondomain{f}{\RR^2}{\RR}\), given by
    \[
        f\pqty{x, y} = x^3 + y^3 - 3xy.
    \]

    We find the gradient as
    \[
        \grad f = \pqty{3 x^2 - 3y, 3 y^2 - 3x},
    \]
    and note
    \[
        \grad f = \vb{0} \iff x^2 = y \text{ and } y^2 = x.
    \]

    This then solves to \(\pqty{x, y} = \pqty{0, 0}\) or \(\pqty{x, y} = \pqty{1, 1}\).

    Now, looking at the Hessian matrix, we have
    \[
        \vb{H} = \begin{pmatrix}
            6x & -3 \\ -3 & 6y
        \end{pmatrix}.
    \]

    The determinant and the trace are therefore given by
    \[
        \det \vb{H} = 9 \pqty{4xy - 1}, \quad \tr \vb{H} = 6 \pqty{x + y}.
    \]

    At \(\pqty{x, y} = \pqty{1, 1}\), \(\det \vb{H} = 27 > 0\) and \(\tr \vb{H} = 12 > 0\). Therefore, \(\pqty{1, 1}\) is a local minimum, and
    \[
        f\pqty{1, 1} = -1.
    \]

    At \(\pqty{x, y} = \pqty{0, 0}\), \(\det \vb{H} = -9 < 0\). Therefore, \(\pqty{0, 0}\) is a saddle, and
    \[
        f\pqty{0, 0} = 0.
    \]

    We may calculate the eigenvalues and the corresponding eigenvectors as
    \[
        \lambda_1 = -3, \quad \vb{e}_1 = \begin{pmatrix}
            1 \\ 1
        \end{pmatrix}
    \]
    and
    \[
        \lambda_2 = 3, \quad \vb{e}_2 = \begin{pmatrix}
            1 \\ -1
        \end{pmatrix}.
    \]

    Therefore, \(\pqty{0, 0}\) is a local maximum along the line parallel to \(\vb{e}_1\) (\(y = x\)), while it is a local minimum along the line parallel to \(\vb{e}_2\) (\(y = -x\)).

    But note \(f\) does not have a global maximum or minimum -- we can make \(x\) and \(y\) very large (but of the same sign), and \(f\) increases/decreases without bound.
\end{example}